\section{Threat Model}
\label{sec:threat}

\textbf{Attack surface.}
The threshold-calibration stage occurs after global model convergence and before threshold-dependent detection begins.
Each client $i$ collects a benign calibration set $\mathcal{C}_i = \{s_1, \ldots, s_{n_i}\}$ of reconstruction scores from its own traffic, then sets its local threshold $\tau_i = \mathrm{quantile}(\mathcal{C}_i,\, q)$ at a fixed quantile $q$ (here $q{=}0.95$).
Depending on the federation's threshold policy, $\tau_i$ may be used directly (\textsc{Local}), averaged across clients (\textsc{Global}), or averaged within a calibration-score cluster (\textsc{Cluster}).

No existing FL training-phase or aggregation-phase defense monitors or validates the calibration buffer: it is collected autonomously by the client after the aggregation round completes.

\textbf{Adversary model.}
The adversary controls a single compromised client $v$, which we call the \emph{poisoned client}.
Under \textsc{Local} policy, only client $v$'s own detection degrades; under \textsc{Global} policy, all federation members receive the contaminated threshold.
We use ``victim'' to refer to client $v$ itself when measuring per-client detection harm, and ``non-victim clients'' when referring to uninvolved federation members.

The adversary has \emph{gray-box score-level access}: since the adversary controls the client, it can read and modify the local calibration buffer before threshold computation.
The adversary does not access model weights, modify gradients, alter other clients' data, or touch test data, test labels, or training labels.
This is a strictly weaker capability than model-poisoning adversaries, which require gradient injection.
The adversary operates at the score level; generating traffic that achieves a target reconstruction score is out of scope (end-to-end realizability is discussed in \S\ref{sec:discussion}).
The \textsc{Replace-Fixed-Budget} mechanism samples with replacement from the tail reservoir, which at high fractions produces duplicate scores detectable by a uniqueness check (disclosed in \S\ref{sec:discussion}).

\textbf{Attack mechanics.}
We evaluate two attack objectives:
\begin{description}
  \item[Raise.] \textsc{Threshold-Raise} (primary): for $f{>}0$, replace $m=\max\!\bigl(1,\mathrm{round}(fn_v)\bigr)$ positions in $\mathcal{C}_v$ with scores drawn with replacement from the top-10\% tail of $\mathcal{C}_v$ (\textsc{High-Score-Benign} source).
  All replacement scores are semantically benign; AUROC is invariant by construction (test scores and labels are unchanged under poisoning).
  A raised threshold $\tau_v$ increases the detection threshold, suppressing alarms at the victim site.
  \item[Lower.] \textsc{Threshold-Lower} (secondary, conditional): replace $m$ positions with scores from the bottom-10\% tail (\textsc{Low-Score-Benign}).
  A lowered threshold increases the false-positive rate (alarm burden) and FPR disparity; the resulting $\Delta\mathrm{TPR}{>}0$ reflects threshold displacement and is a secondary indicator of operating-point shift.
  This attack is structurally conditional: random position replacement displaces upper-tail support with probability $\approx f$.
\end{description}

\textbf{Negative control.} \textsc{Random-Benign}: replacement scores are sampled uniformly from the full calibration distribution rather than a tail.
Any effect attributable to score contamination rather than position perturbation should vanish in this condition.

\textbf{Scope boundary.}
This paper characterizes the vulnerability at the score level.
A complete end-to-end attack would additionally require generating network traffic whose autoencoder reconstruction score lands in the target tail---this is future work (\S\ref{sec:discussion}).
The score-level characterization is a controlled experiment, not a deployed attack: it isolates a necessary calibration-channel vulnerability mechanism, not a complete traffic-generation attack.
Claims are bounded to the single-client, single-compromise scenario; multi-client collusion is not evaluated.
